\section{Conclusions}
%

We have presented an analysis of power-suppressed (higher-twist) contributions to qPDFs and pPDFs based on the study of factorial divergences (renormalons) in the corresponding coefficient functions within the bubble-chain approximation. Factorial asymptotic implies that the sum of the series is only defined to a power accuracy and therefore, the QCD perturbation theory must be corrected by nonperturbative power-suppressed contributions to produce unambiguous predictions. Our results have to be considered as a ``minimal model'' for the higher-twist corrections that captures effects that are necessary for the selfconsistency of the theory, but possibly misses other nonperturbative corrections that are, e.g, protected 
by symmetries and not ``seen'' through perturbative expansions. Our main conclusions are as follows:
\begin{itemize}
\item Position space PDFs (qITDs) have flat power corrections at large Ioffe times. Generally, power corrections are much 
      larger for the ``transverse'' projection as compared to the ``longitudinal'' projection. 
\item Power corrections for qPDFs have a generic behavior
\begin{align}
  \cQ (x,p) &=  q(x) \biggl\{1 + \mathcal{O}\left( \frac{\Lambda^2}{p^2}\cdot\frac{1}{{x^2(1-x)}}\right) \biggr\}.     
\end{align}
Note that the corresponding target mass corrections, \eqref{target-parallel} and \eqref{target-perp}, do not show up the  $1/x^2$ enhancement. This behavior is commensurate to the suppression of Nachtmann's target mass corrections $\sim x^2 m^2/Q^2$ at small $x$ for the DIS structure functions. The normalization of the underlying qITDs to unity at the zero momentum considerably reduces the size the power correction to the qPDFs at $ x \lesssim 0.5$ at the cost of a strong enhancement at $x\gtrsim 0.5-0.6$. Thus, such a normalization procedure is not suitable to access the large-$x$ behavior of the PDFs, but apparently minimizes power corrections in the intermediate $x$-region. Note that the above discussion is based on a normalization procedure different from the nonperturbative renormalization used in ref.~\cite{Chen:2018xof}, thus the power corrections for the latter might behave differently from what has been shown here.
 \item
Power corrections for pPDFs have a generic behavior
\begin{align}
\mathcal{P}(x,z) &= q(x)\Big\{1 + \mathcal{O}\big( z^2\!\Lambda^2{(1-x)}\big)\Big\}.
\end{align}
The suppression at $x\to 1$ is lifted by the zero-momentum normalization factor. However, it can be upheld by the normalization to the vacuum matrix element of the same operator. We conclude that pPDFs can offer an interesting alternative to qPDFs for the study of large-$x$ region.
\end{itemize}
As a byproduct of this study, we provide the results for the coefficient functions 
in the large-$n_f$ approximation, terms $\sim n_f^k\alpha_s^{k+1}$,  to all orders in perturbation theory.
These results can be useful to estimate the effects of uncalculated higher orders and scale-setting 
using the BLM-type procedure. The corresponding analysis goes beyond the scope of this paper and will be presented in a future publication.
 
